\documentclass[11pt]{article}

\usepackage[T1]{fontenc}
\usepackage[utf8]{inputenc}
\usepackage{lmodern}
\usepackage{amsmath,amssymb,amsthm,mathtools}
\usepackage{geometry}
\usepackage{booktabs}
\usepackage{array}
\usepackage{longtable}
\usepackage{hyperref}
\usepackage{enumitem}

\geometry{margin=1in}
\setlist[itemize]{leftmargin=1.5em}
\setlist[enumerate]{leftmargin=1.75em}

\newtheorem{remark}{Remark}
\newtheorem{observation}{Observation}

\title{Memorandum on the EasyCrypt Formalization of the HAETAE NTT}
\author{Codex Proof Report}
\date{22 April 2026}

\begin{document}
\maketitle

\begin{abstract}
This note records the present state of the EasyCrypt proof development for the
HAETAE number-theoretic transform, with particular attention to the
refinement path from the extracted Jasmin program to the abstract polynomial
specification. The report is written in the style of a mathematical proof
memorandum: completed lemmas are stated by role, partially reconstructed
arguments are isolated precisely, and the remaining obligations are organized
according to the logic of the final correctness theorem. The principal
technical achievements to date are the proof of losslessness for the
imperative NTT specifications, the reconstruction of the root-of-unity
identity $zroot^{256} = -1$ in $\mathbf{F}_q$, and the establishment of a
substantial portion of the representation layer relating machine words and
mathematical coefficients. The main unresolved point at present is a
library-level normalization issue in the concrete twiddle-table bridge,
specifically around the rewriting of \texttt{BArray1024.get32\_of\_list32}
inside the full theory \texttt{NTT\_Fq.ec}.
\end{abstract}

\section{Objective and Verification Architecture}

The final theorem sought in the project is an end-to-end refinement statement:
the Jasmin routines implementing the forward and inverse NTT must coincide,
under a suitable representation relation, with the mathematical
transformations defined over the quotient ring formalized in
\texttt{Rq.ec}. In conceptual terms, the proof must traverse four layers:
\[
\text{Jasmin source}
\;\Longrightarrow\;
\text{EasyCrypt extraction}
\;\Longrightarrow\;
\text{imperative EasyCrypt NTT model}
\;\Longrightarrow\;
\text{abstract ring-theoretic NTT specification}.
\]

The relevant artifacts are as follows.
\begin{itemize}
\item \texttt{extract/Hpoly\_extract.ec}: extracted Jasmin code, including the concrete tables \texttt{jzetas} and \texttt{jzetas\_inv}.
\item \texttt{NTT\_Fq.ec}: imperative NTT specifications together with the algebraic and representation lemmas required to relate program state to mathematical state.
\item \texttt{Fq.ec} and \texttt{Montgomery.ec}: arithmetic correctness of modular and Montgomery operations.
\item \texttt{Rq.ec}: abstract forward and inverse NTT operators.
\end{itemize}

\section{Foundational Proofs Already Established}

\subsection{Losslessness of the Imperative Specifications}

The imperative module \texttt{NTT} in \texttt{NTT\_Fq.ec} contains two
procedures:
\begin{itemize}
\item \texttt{NTT.ntt}, modeling the forward transform with eight butterfly layers;
\item \texttt{NTT.invntt}, modeling the inverse transform with eight inverse layers followed by the final scaling by $\operatorname{inv}(256)$.
\end{itemize}

Both procedures are now accompanied by complete losslessness proofs:
\begin{align*}
&\texttt{lemma ntt\_spec\_ll : islossless NTT.ntt},\\
&\texttt{lemma invntt\_spec\_ll : islossless NTT.invntt}.
\end{align*}

These proofs proceed by explicit integer variants for the three nested loops.
For the forward transform one uses:
\begin{itemize}
\item outer variant: \texttt{len};
\item middle variant: $256 - \texttt{start}$;
\item inner variant: $\texttt{start} + \texttt{len} - \texttt{j}$.
\end{itemize}
For the inverse transform the final scaling loop contributes the additional
variant $256 - \texttt{j}$, while the outer loop uses $256 - \texttt{len}$.

The supporting lemmas
\[
\texttt{ntt\_middle\_right},\qquad
\texttt{inv\_middle\_right},\qquad
\texttt{inv\_outer\_progress},\qquad
\texttt{ntt\_init\_progress}
\]
encode the required arithmetic monotonicity of these variants. Their role is
entirely proof-theoretic: they guarantee that the imperative models are
legitimate targets for subsequent pRHL and Hoare refinements.

\subsection{Root-of-Unity Arithmetic}

The HAETAE parameters are fixed by
\[
q = 64513,\qquad N = 256,\qquad zroot = 426 \in \mathbf{F}_q.
\]
The crucial algebraic fact needed for the inverse transform is now proved:
\[
\texttt{lemma exp\_zroot\_256 : Zq.exp zroot 256 = incoeff (-1)}.
\]

The argument is economical. One first identifies
\[
zroot^{128} = 28837
\]
by reading the value at index \(1\) in the explicit forward zeta table, and
then derives
\[
zroot^{256} = (zroot^{128})^2 = -1 \pmod q.
\]
This confirms that \(zroot\) behaves as a primitive \(512\)-th root of unity,
which is precisely the structural fact needed to justify the inverse
normalization.

\section{Representation Layer Between Machine Words and Coefficients}

The extracted Jasmin code operates on byte arrays of packed \(32\)-bit words,
whereas the mathematical specification is expressed over arrays of field
coefficients. The present development already contains the nucleus of the
representation relation:
\begin{align*}
&\texttt{word\_to\_coeff}(w) = \texttt{incoeff}(\texttt{W32.to\_sint}\; w),\\
&\texttt{barray256\_to\_poly}(rp) =
  \texttt{Array256.init}(\lambda i.\,\texttt{word\_to\_coeff}(\texttt{BArray1024.get32}\; rp\; i)),\\
&\texttt{poly\_repr}(rp,p) \equiv \texttt{barray256\_to\_poly}(rp) = p.
\end{align*}

This layer is supported by several useful conversion lemmas:
\begin{itemize}
\item \texttt{barray256\_to\_polyE};
\item \texttt{poly\_repr\_get};
\item \texttt{word\_to\_coeff\_of\_int};
\item \texttt{word\_to\_coeff\_small};
\item \texttt{word\_to\_coeff\_repr};
\item the concrete instances \texttt{word\_to\_coeff\_neg16505} and \texttt{word\_to\_coeff\_26964}.
\end{itemize}

Mathematically, these lemmas do the indispensable bookkeeping which allows one
to pass from signed machine integers to residue classes in \(\mathbf{F}_q\).
Cryptographically, this is the bridge by which a concrete read from memory may
eventually be interpreted as a field element participating in a butterfly.

\section{Explicit Twiddle Tables and Their Formal Role}

\subsection{Inverse Twiddle Table}

The inverse zeta table is defined abstractly by
\[
\texttt{zetas\_inv}[k] = zroot^{-(\operatorname{bsrev}_8(k)+1)}
\quad (0 \le k \le 254),
\]
with the final entry replaced by
\[
\texttt{scale255} = \operatorname{inv}(256) = -252 \pmod q.
\]

The following items are now fully formalized in \texttt{NTT\_Fq.ec}:
\begin{itemize}
\item \texttt{unit\_zroot};
\item \texttt{inv\_zroot};
\item \texttt{zetas\_invE}, giving the full explicit list of inverse twiddle factors;
\item \texttt{scale255E};
\item \texttt{zetas\_inv\_vals}, expressing the Montgomery-scaled inverse table explicitly;
\item \texttt{zetas\_inv\_255\_montE};
\item \texttt{jzetas\_inv\_255\_wordE};
\item \texttt{jzetas\_inv\_255\_scaleE}.
\end{itemize}

The last two lemmas are especially significant. The extracted table
\texttt{jzetas\_inv} stores at index \(255\) not merely the abstract factor
\(\operatorname{inv}(256)\), but the machine-level constant
\(-29720\), representing the Montgomery-encoded quantity required by the
Jasmin multiplication routine. The lemma
\[
\texttt{jzetas\_inv\_255\_scaleE}
\]
establishes the exact compatibility relation
\[
\texttt{word\_to\_coeff}(\texttt{jzetas\_inv}[255]) \cdot R^{-1}
=
\texttt{scale255} \cdot R,
\]
which is precisely the identity needed in the final scaling loop of the
inverse transform.

\subsection{Forward Twiddle Table}

The forward table is treated similarly. The abstract definition
\[
\texttt{zetas}[k] = zroot^{\operatorname{bsrev}_8(k)}
\]
is matched by the explicit formulas
\begin{itemize}
\item \texttt{zetasE};
\item \texttt{zetas\_vals}.
\end{itemize}

Moreover, the development already contains the index-\(1\) bridge lemmas
\begin{itemize}
\item \texttt{jzetas\_1\_poly};
\item \texttt{jzetas\_1\_word\_montE};
\item \texttt{jzetas\_1\_get32E};
\item \texttt{jzetas\_1\_get32dE};
\item \texttt{zetas\_1\_montE};
\item \texttt{jzetas\_1E};
\item \texttt{jzetas\_1\_wordE}.
\end{itemize}

These serve as small but concrete witnesses that the forward table relation is
mathematically the expected one. They also indicate the intended pattern for
the full array-level proof.

\section{Partial Reconstruction of the Twiddle-Bridge Proofs}

The principal work currently in progress concerns the array-level bridge
lemmas relating the extracted tables \texttt{jzetas} and
\texttt{jzetas\_inv} to the mathematical arrays
\texttt{array256\_mont zetas} and \texttt{array256\_mont\_inv zetas\_inv}.

The present proof scripts introduce the stronger intermediate identities
\begin{align*}
&\texttt{jzetas\_inv\_polyE}:
  \texttt{barray256\_to\_poly Hpoly\_extract.jzetas\_inv}
  =
  (\texttt{array256\_mont\_inv zetas\_inv})[255 \leftarrow \texttt{incoeff}(-29720)],\\
&\texttt{jzetas\_polyE}:
  \texttt{barray256\_to\_poly Hpoly\_extract.jzetas}
  =
  (\texttt{array256\_mont zetas})[0 \leftarrow \texttt{incoeff}(0)].
\end{align*}

From these one deduces the pointwise statements
\begin{align*}
&\texttt{jzetas\_inv\_poly\_vals}\; i
  \quad (0 \le i < 255),\\
&\texttt{jzetas\_poly\_vals}\; i
  \quad (1 \le i < 256),
\end{align*}
and thence the direct get-lemmas \texttt{jzetas\_inv\_get} and
\texttt{jzetas\_get}.

\begin{observation}
The special indices differ for conceptual reasons.
For the inverse table, index \(255\) stores the Montgomery-encoded final
normalization constant, not the plain array value
\((\texttt{array256\_mont\_inv zetas\_inv})[255]\).
For the forward table, index \(0\) is padding and is never read by the
algorithm, because the counter is incremented before the first lookup.
\end{observation}

\subsection{Current Obstruction}

The mathematical structure of the proof is clear and essentially complete.
The obstruction is local and syntactic. When one compiles
\texttt{NTT\_Fq.ec}, EasyCrypt currently fails in the proof of
\texttt{jzetas\_inv\_polyE} at the line
\[
\texttt{rewrite BArray1024.get32\_of\_list32 1:/\#}.
\]
The prover reports:
\begin{quote}
\texttt{nothing to rewrite}.
\end{quote}

Thus the unresolved issue is not arithmetic but normalization of the concrete
\texttt{BArray1024} accessor inside the full ambient theory. In isolated probe
files the same rewrite shape is admissible, whereas inside the complete theory
it fails to match the current goal term. This suggests a mismatch in the
surface form of the accessor after unfolding \texttt{Hpoly\_extract.jzetas}
and \texttt{Hpoly\_extract.jzetas\_inv}, rather than any defect in the
underlying mathematical argument.

\section{Status of the Reconstruction}

The present state of the formal development may be summarized as follows.

\begin{center}
\begin{longtable}{p{0.27\linewidth} p{0.21\linewidth} p{0.43\linewidth}}
\toprule
Component & Status & Mathematical significance \\
\midrule
\endfirsthead
\toprule
Component & Status & Mathematical significance \\
\midrule
\endhead
\texttt{ntt\_spec\_ll} & proved & termination and losslessness of the forward imperative model \\
\texttt{invntt\_spec\_ll} & proved & termination and losslessness of the inverse imperative model \\
\texttt{exp\_zroot\_256} & proved & primitive-root structure and inverse scaling justification \\
\texttt{word\_to\_coeff} layer & proved & conversion from signed words to field coefficients \\
\texttt{poly\_repr} layer & proved & base representation predicate for future simulation arguments \\
\texttt{zetas\_invE}, \texttt{zetas\_inv\_vals} & proved & explicit inverse twiddle constants and Montgomery form \\
\texttt{zetasE}, \texttt{zetas\_vals} & proved & explicit forward twiddle constants and Montgomery form \\
\texttt{jzetas\_inv\_255\_scaleE} & proved & compatibility of the inverse final scaling constant \\
\texttt{jzetas\_inv\_polyE} and descendants & formalized but not compiling & inverse array bridge blocked by accessor normalization \\
\texttt{jzetas\_polyE} and descendants & formalized but not compiling & forward array bridge blocked by the same normalization issue \\
\texttt{fqmul\_corr} & not yet formalized & arithmetic correctness of one Montgomery multiplication \\
Butterfly correctness lemmas & not yet formalized & local simulation step for both transforms \\
Loop invariants for extracted Jasmin code & not yet formalized & pRHL refinement from extraction to imperative model \\
Imperative-to-abstract refinement & not yet formalized & identification of imperative NTT with \texttt{Rq.ntt} and \texttt{Rq.invntt} \\
\bottomrule
\end{longtable}
\end{center}

\section{Remaining Proof Obligations}

Once the concrete table bridge is resolved, the remaining development is
structurally straightforward, though not short.

\subsection{Arithmetic Step}

One must prove correctness of the extracted routine
\texttt{M.\_\_fqmul} using Montgomery reduction, relying on the
existing theory \texttt{Fq.ec} and the lemma \texttt{SREDCp\_corr} from
\texttt{Montgomery.ec}. This yields the field-level interpretation of one
machine multiplication.

\subsection{Local Butterfly Lemmas}

With the twiddle tables and multiplication routine in place, the next step is
to establish:
\begin{itemize}
\item a forward butterfly correctness lemma;
\item an inverse butterfly correctness lemma, including the final scaling
phase.
\end{itemize}

These are the atomic transitions from which the loop invariants for the
extracted procedures will be assembled.

\subsection{Loop Refinement}

The extracted routines \texttt{M.\_poly\_ntt} and \texttt{M.\_poly\_invntt}
must then be related to the imperative procedures \texttt{NTT.ntt} and
\texttt{NTT.invntt} by synchronized invariants on the indices
\texttt{len}, \texttt{start}, \texttt{j}, and \texttt{zetasctr}. The natural
logic here is probabilistic relational Hoare logic, though the programs are
deterministic.

\subsection{Abstract Identification}

Finally, one must prove that the imperative models coincide with the abstract
ring-level operators
\[
\texttt{Rq.ntt}, \qquad \texttt{Rq.invntt}.
\]
This is where the Cooley--Tukey decomposition and the bit-reversal structure
must be identified with the explicit formula from \texttt{Rq.ec}.

\section{Recommended Next Step}

The proof search should now focus narrowly on the accessor normalization
problem in the twiddle-bridge lemmas. Concretely:
\begin{enumerate}
\item inspect the exact normal form of
  \texttt{BArray1024.get32 Hpoly\_extract.jzetas\_inv i}
  after unfolding \texttt{Hpoly\_extract.jzetas\_inv} inside the ambient theory;
\item identify the precise theorem name or rewriting direction accepted by the
  cloned \texttt{ByteArray} theory in this context;
\item only after that, resume the array extensional proofs of
  \texttt{jzetas\_inv\_polyE} and \texttt{jzetas\_polyE}.
\end{enumerate}

\begin{remark}
At the present moment the main body of mathematical work is not blocked by a
deep algebraic deficiency. The obstruction is a proof-engineering issue in the
representation layer. This is good news: once the concrete memory accessor is
brought under control, the path to the local butterfly lemmas is already
visible.
\end{remark}

\section*{Conclusion}

The proof effort has already secured the algebraic spine of the HAETAE NTT
formalization. The imperative specifications are lossless; the chosen root is
shown to have the required order; the explicit zeta tables, both forward and
inverse, are mathematically identified; and the representation relation from
word arrays to coefficient arrays is in place in its essential form. What
remains is the completion of the concrete twiddle-table bridge and, beyond it,
the expected sequence of simulation lemmas from Montgomery multiplication to
butterflies, from butterflies to loops, and from loops to the abstract NTT
operators. The present state is therefore best described not as incomplete in
principle, but as suspended at a technically local juncture.

\end{document}
